% C086 | 双曲线焦距比与离心率取值范围（简短源稿；不依赖图）
% 供主模板 \input 的 LaTeX 片段，非独立文档。
\textbf{C086\quad 双曲线焦距比与离心率取值范围}

设双曲线 $H:\dfrac{x^2}{a^2}-\dfrac{y^2}{b^2}=1$（$a,b>0$），
焦点 $F_1(-c,0)$、$F_2(c,0)$，$c^2=a^2+b^2$，$e=c/a>1$。
给定 $\lambda>0$，双曲线上存在点 $P$ 满足
$PF_1=\lambda PF_2$ 的\textbf{充要条件}为
\[
\boxed{\lambda\ne1,\qquad 1<e\le\frac{\lambda+1}{|\lambda-1|}}.
\]
\textbf{最快证明：}令 $r_i=PF_i$。双曲线定义给出
$|r_1-r_2|=2a$，三角形不等式给出 $r_1+r_2\ge 2c$；
因此 $e=c/a\le (r_1+r_2)/|r_1-r_2|=(\lambda+1)/|\lambda-1|$。
反之，按比例取 $r_1=2a\lambda/|\lambda-1|$、
$r_2=2a/|\lambda-1|$，则 $|r_1-r_2|=2a<2c\le r_1+r_2$；
两焦点为圆心的两个圆存在交点 $P$，且该交点必在双曲线上。
\textbf{几何本质：}定比轨迹 $PF_1=\lambda PF_2$（$\lambda\ne1$）
是阿波罗尼斯圆，圆心 $C=(c(\lambda^2+1)/(\lambda^2-1),0)$、
半径 $R=2c\lambda/|\lambda^2-1|$；它与\emph{定差轨迹}双曲线是否相交，
决定有无所求点。
\textbf{等号：}$P$ 恰好位于对应支的顶点；严格小于时有关于 $x$ 轴对称的两点。
